% !TeX program = lualatex
% !TeX encoding = utf8
% !TeX spellcheck = uk_UA
% !TeX root =../PyscfBook.tex

%% ============================================================
\chapter{Молекулярні властивості як похідні енергії}
%% ============================================================

У квантовій хімії енергія $E$ молекули є основним об'єктом розрахунків, але
вона не єдина величина, що визначає поведінку системи.
Переважна більшість експериментальних спостережуваних --- це \emph{похідні
	енергії} за певними зовнішніми параметрами:
електричним полем, магнітним полем, координатами ядер, формою потенціалу,
навантаженням тощо.
Саме тому аналітичні похідні енергії відіграють фундаментальну роль
у квантовій теорії молекулярних властивостей.

\section{Молекулярні властивості як похідні енергії}

Енергія квантової системи залежить від зовнішнього параметра $x$,
який може описувати дію електричного або магнітного поля,
зміну геометрії молекули, деформацію потенціалу тощо:
\begin{equation}
	E = E(x).
	\label{eq:E_x_def}
\end{equation}

Для малих змін параметра виконують ряд Тейлора:
\begin{equation}
	E(x)=
	E^{(0)} + E^{(1)}x + \tfrac12 E^{(2)}x^2 + \tfrac16 E^{(3)}x^3 + \dots,
	\label{eq:E_Taylor}
\end{equation}
де
\begin{equation}
	E^{(n)} =
	\left.\frac{d^n E}{dx^n}\right|_{x=0}.
	\label{eq:E_n_def}
\end{equation}

Кожна похідна має фізичний зміст, оскільки описує реакцію енергії на зовнішнє збурення. Розглянемо кілька фундаментальних прикладів.

\paragraph{Дипольний момент та поляризовність.}

Почнемо з фундаментального визначення.
Дипольний момент молекули --- це зміна енергії при малому збільшенні
однорідного електричного поля $\mathcal{E}$:
\begin{equation}
	\mu = -\,\left.\frac{dE}{d\mathcal{E}}\right|_{\mathcal{E}=0}.
	\label{eq:dipole_def}
\end{equation}

Іншими словами, дипольний момент — це лінійний коефіцієнт у розкладі
енергії за полем. Однак цю ж думку можна подати з іншого боку:
розглянемо розклад самого дипольного моменту за полем.

Дипольний момент у полі $\mathcal{E}$ змінюється так:
\begin{equation}
	\mu(\mathcal{E})
	=
	\mu^{(0)}
	+
	\alpha\,\mathcal{E}
	+
	\tfrac12 \beta\,\mathcal{E}^2
	+\dots ,
	\label{eq:mu_Taylor}
\end{equation}
де
\begin{itemize}
	\item $\mu^{(0)}$ --- власний дипольний момент у відсутності поля;
	\item $\alpha$ --- поляризованість;
	\item $\beta$ --- друга гіперполяризованість;
	\item наступні члени відповідають нелінійному відгуку.
\end{itemize}

\textbf{Загалом}, повний дипольний момент \(\mu(\mathcal{E})\) складається з:
\begin{itemize}
\item \textbf{перманентного дипольного моменту} \(\mu^{(0)}\) — існує завжди (якщо \(\mu^{(0)} \neq 0\)),
\item \textbf{наведеного дипольного моменту} — виникає під дією поля і включає як лінійний (\(\alpha \mathcal{E}\)), так і нелінійні (\(\beta \mathcal{E}^2\), \(\gamma \mathcal{E}^3\) тощо) вклади.
\end{itemize}

Тепер використаємо факт, що енергія --- це інтеграл взаємодії дипольного моменту з полем:
\begin{equation}
	E(\mathcal{E})
	=
	E_0 - \int_0^{\mathcal{E}} \mu(\mathcal{E}')\, d\mathcal{E}'.
	\label{eq:E_as_integral_of_mu}
\end{equation}

Підставляємо розклад \eqref{eq:mu_Taylor}:
\begin{equation}
	E(\mathcal{E}) =
	E_0
	- \mu^{(0)}\mathcal{E}
	- \tfrac12 \alpha\,\mathcal{E}^2
	- \tfrac16 \beta\,\mathcal{E}^3
	- \dots
	\label{eq:E_expanded_correct}
\end{equation}

З цього рівняння вже видно, що:
\begin{equation}
	\mu^{(0)}  = -\,\left.\frac{dE}{d\mathcal{E}}\right|_{\mathcal{E}=0}, \quad
	\alpha     = \left.\frac{d^2E}{d\mathcal{E}^2}\right|_{\mathcal{E}=0}, \quad
	\beta      = -\,\left.\frac{d^3E}{d\mathcal{E}^3}\right|_{\mathcal{E}=0},
\end{equation}

тобто:
\[
	\text{перша похідна енергії дає дипольний момент,}
\]
\[
	\text{друга похідна енергії дає поляризованість,}
\]
\[
	\text{третя похідна --- гіперполяризованість.}
\]

Саме тому всі електричні властивості природно визначаються
похідними енергії за зовнішнім полем.

\paragraph{Сили на ядрах.}

Почнемо з фундаментального визначення.
Сила, що діє на ядро $A$, визначається як зміна енергії при малому зміщенні
координати $R_A$:
\begin{equation}
	F_A = -\,\frac{dE}{dR_A}.
	\label{eq:force_def}
\end{equation}

Це є пряме квантово-механічне формулювання принципу:
\[
	\text{сила} = -\nabla E.
\]

Розглянемо тепер, як сила змінюється при малих зміщеннях ядра.
Вона також може бути розкладена в ряд Тейлора:
\begin{equation}
	F_A(R_A)
	=
	F_A^{(0)}
	+
	\sum_B k_{AB}\,(R_B - R_B^{(0)})
	+
	\dots ,
	\label{eq:force_Taylor}
\end{equation}
де $k_{AB}$ --- коефіцієнти, що описують зміну сили при зміщенні інших ядер.

Коефіцієнти $k_{AB}$ --- це елементи матриці силових констант:
\begin{equation}
	k_{AB} = \left.\frac{\partial F_A}{\partial R_B}\right|_{0}.
	\label{eq:force_constants_def}
\end{equation}

Тепер скористаємося зв’язком між силою і енергією \eqref{eq:force_def}.
Диференціюємо його за $R_B$:
\begin{equation}
	\frac{\partial F_A}{\partial R_B}
	=
	-\,\frac{\partial^2 E}{\partial R_A \partial R_B}.
	\label{eq:F_derivative_relation}
\end{equation}

Таким чином,
\[
	k_{AB} = -\,\frac{\partial^2 E}{\partial R_A \partial R_B}.
\]

З іншого боку, друга похідна енергії природно визначає елемент матриці гессіана:
\begin{equation}
	H_{AB} =
	\frac{\partial^2 E}{\partial R_A\,\partial R_B}.
	\label{eq:nuclear_Hessian_def_fixed}
\end{equation}

Отже,
\[
	k_{AB} = -H_{AB}.
\]

\paragraph{Ядерний гесiян.}

Ядерний гесiян --- це другий диференціал енергії за координатами ядер:
\begin{equation}
	H_{AB} =
	\frac{\partial^2 E}{\partial R_A\,\partial R_B},
\end{equation}
і саме ця матриця визначає:

\begin{itemize}
	\item форму та кривизну поверхні потенціальної енергії,
	\item рівноважність геометрії (мінімум, сідлові точки),
	\item вібраційні частоти та нормальні координати.
\end{itemize}

Зв’язок
\[
	F_A = -\frac{\partial E}{\partial R_A}, \qquad
	k_{AB} = \frac{\partial F_A}{\partial R_B},
\]
призводить до висновку:
\[
	H_{AB} = -\,k_{AB}.
\]

Таким чином, як електричні,так і механічні властивості молекули
описуються єдиним математичним підходом --- через похідні енергії за параметрами, що збурюють систему.

Іншими словами, похідні енергії визначають, як молекула реагує на поля або зміни геометрії. Цю  реакцію системи на зовнішні збурення називають відгуком.

Концепція відгуку надзвичайно універсальна: практично всі молекулярні властивості можна розглядати як похідні енергії за зовнішніми параметрами --- електричним або магнітним полем, координатами ядер (саме цей випадок був детально розглянутий в розділі \ref{Oscillations}) тощо.
В табл. \ref{tab:prop}) наведено узагальнену схему, що показує, як похідні різних порядків визначають ті чи інші властивості .

\begin{commentbox}
\textbf{Теорія відгуку} --- це узагальнення теорії збурень, в якій
замість формального розкладу хвильової функції розглядається її \emph{варіаційний відгук} на зовнішнє поле.
\end{commentbox}




\begin{table}[h!]
\centering
\caption{Зв’язок порядку похідних енергії з молекулярними властивостями. Тут $q$ --- координати, $E$ --- зовнішнє електричне поле, $B$ --- зовнішнє магнітне поле, $I$ --- магнітне поле ядра.}\label{tab:prop}
\begin{tblr}{
%  cells = {font=\scriptsize},
  colspec={Q[c,m ] Q[c,m] Q[c,m] Q[c,m] Q[l,m]},
  cell{1}{1} = {c=4}{c},
  cell{1}{5} = {}{c,m},
  cell{2-Z}{1-4}={}{cmd=\bfseries},
  stretch = -1,
  hline{2} = {2pt},
  hline{3} = {1pt},
  hline{Z} = {1pt},
}
$n$-та похідна енергії &       &       &       & {Збурення} \\
$n_R$ & $n_E$ & $n_B$ & $n_I$ & Властивість \\
1 & 0 & 0 & 0 & Сила \\
2 & 0 & 0 & 0 & Частоти гармонічних коливань \\
3 & 0 & 0 & 0 & Ангармонічні поправки \\
0 & 1 & 0 & 0 & Електричний дипольний момент $\vec{\mu}_e$ \\
0 & 2 & 0 & 0 & Електрична поляризованість $\chi_e$ \\
0 & 3 & 0 & 0 & Електрична надполяризованість \\
0 & 0 & 1 & 0 & Магнітна сприйнятливість $\chi_m$ \\
0 & 0 & 2 & 1 & Константа надтонкого зв’язку (\texttt{EPR}) \\
0 & 0 & 0 & 2 & Стала спін-спінової взаємодії ядер \\
1 & 1 & 0 & 0 & Інтенсивності ІЧ-спектрів \\
2 & 1 & 0 & 0 & Інтенсивності обертонів в ІЧ-спектрах \\
1 & 2 & 0 & 0 & Інтенсивності Раманівських спектрів \\
2 & 2 & 0 & 0 & Інтенсивності обертонів у Раманівських спектрах \\
0 & 1 & 1 & 0 & Циркулярний дихроїзм \\
0 & 2 & 1 & 0 & Магнітний циркулярний дихроїзм \\
0 & 0 & 1 & 1 & Хімічний зсув (\texttt{NMR}) \\
\end{tblr}
\end{table}



Як видно, усі ці властивості --- від сил та частот до магнітного екранування ---
описуються похідними енергії різних порядків за різними збуреннями.
Таким чином, \emph{усі властивості молекули є лише різними проявами одного й того ж принципу: відгуку хвильової функції на зовнішнє збурення.}

%% --------------------------------------------------------
\section{Чисельні та  аналітичні похідні}
%% --------------------------------------------------------

Найпростіший спосіб отримати похідну енергії --- використати кінцеві різниці.

\paragraph{Перша похідна.}

Одностороння формула (перший порядок точності):
\begin{equation}
	\frac{dE}{dx}
	\approx
	\frac{E(x+h)-E(x)}{h}.
	\label{eq:finite_diff_forward}
\end{equation}

Центральна формула (другий порядок точності):
\begin{equation}
	\frac{dE}{dx}
	\approx
	\frac{E(x+h)-E(x-h)}{2h}.
	\label{eq:finite_diff_central_1st}
\end{equation}

\paragraph{Друга похідна.}

Центральна формула для другої похідної:
\begin{equation}
	\frac{d^2E}{dx^2}
	\approx
	\frac{E(x+h)-2E(x)+E(x-h)}{h^2}.
	\label{eq:finite_diff_2nd}
\end{equation}

\paragraph{Змішана друга похідна.}

Для двох параметрів $x$ і $y$ змішана похідна наближається як
\begin{equation}
	\frac{\partial^2 E}{\partial x\,\partial y}
	\approx
	\frac{
		E(x+h,\,y+k)
		-
		E(x+h,\,y-k)
		-
		E(x-h,\,y+k)
		+
		E(x-h,\,y-k)
	}{4hk}.
	\label{eq:finite_diff_mixed}
\end{equation}

Однак такий підхід має суттєві недоліки:

\begin{itemize}
	\item похибка різницевого диференціювання, яка зростає при малих кроках;
	\item необхідність багатьох повних SCF-розрахунків;
	\item неможливість отримати точні частоти коливань у великих системах;
	\item нестабільність для високих похідних та слабких збурень.
\end{itemize}

У той час як \emph{аналітичні} формули забезпечують:
\begin{itemize}
	\item точні градієнти і гессіани без чисельних похибок;
	\item значно нижчу обчислювальну складність;
	\item природне включення релаксації електронів;
	\item можливість обробляти великі молекули та складні властивості.
\end{itemize}

Саме тому в сучасній квантовій хімії більшість властивостей ---
від геометрій до поляризованостей, силових констант і спектрів ---
за можливості обчислюються через аналітичні похідні.
Аналітичні методи мають високу точність і масштабованість,
оскільки дозволяють уникнути чисельних кінцевих різниць.

Втім, аналітичні формули доступні не для всіх методів і не для всіх
властивостей. У складніших теоріях (наприклад, \texttt{CCSD(T)} або
деяких багатоконфігураційних підходах) повні аналітичні похідні можуть
бути недоступні або надто дорогі.
У таких випадках застосовують обмежені аналітичні похідні, гібридні
схеми або чисто чисельні оцінки.


Важливо відзначити, що аналітичні похідні можливі лише тому, що
енергія $E$ у методах \texttt{HF} та \texttt{DFT} має
замкнений аналітичний вигляд.
Енергія виражається через інтеграли від відомих базисних функцій,
а всі оператори (кулонівський, обмінний, зовнішнього поля) мають
чітко визначені математичні формули.
Тому похідні будь-якого порядку можна отримати шляхом
аналітичного диференціювання цих виразів.

Іншими словами,
наявність явного аналітичного виразу для енергії дозволяє
обчислювати $\frac{dE}{dx}$, $\frac{d^2E}{dx^2}$,
а за потреби й похідні вищих порядків,
не вдаючись до чисельних кінцевих різниць.


\section{Варіаційні методи та залежність енергії від параметрів}

Багато важливих квантовохімічних методів, зокрема
\texttt{HF}, \texttt{DFT} і \texttt{CASSCF}, є варіаційними.
Це означає, що їхня хвильова функція $\Psi$ залежить від набору
варіаційних параметрів $\mathbf p$, які визначають молекулярні орбіталі.
Енергію таких методів можна подати у вигляді
\begin{equation}\label{eq:E_x_p}
E = E(x, \mathbf p(x)),
\end{equation}
де $x$ --- зовнішній параметр (координата ядра, електричне поле,
зміна геометрії тощо).

Позначимо через \(\mathbf p_0(x)\) оптимальні (стаціонарні) параметри для даного \(x\),
тобто ті, що задовольняють умову стаціонарності:
\begin{equation}
\left.\frac{\partial E(x,\mathbf p)}{\partial p_i}\right|_{\mathbf p=\mathbf p_0(x)} = 0
\qquad\text{для всіх }i.
\label{eq:stationarity_p0}
\end{equation}
У варіаційних методах (\texttt{HF}, \texttt{DFT}, \texttt{CASSCF})
ця умова виконується умова.


Тоді повна похідна енергії при зміні \(x\) записується як:
\begin{equation}
\frac{dE}{dx}
=
\left.\frac{\partial E(x,\mathbf p)}{\partial x}\right|_{\mathbf p=\mathbf p^{(0)}(x)}
+
\sum_i
\left.\frac{\partial E(x,\mathbf p)}{\partial p_i}\right|_{\mathbf p=\mathbf p^{(0)}(x)}
\frac{d p_{0i}(x)}{dx}.
\label{eq:dE_dx_general}
\end{equation}

Оскільки для оптимальних параметрів \(\mathbf p^{(0)}(x)\) виконано \eqref{eq:stationarity_p0},
другий член дорівнює нулю, і отримуємо робочу формулу для першої похідної:
\begin{equation}
\frac{dE}{dx}
=
\left.\frac{\partial E(x,\mathbf p)}{\partial x}\right|_{\mathbf p=\mathbf p^{(0)}(x)}.
\label{eq:dE_dx_stationary}
\end{equation}
Видно, що \emph{перші} похідні залежать лише від явної залежності $E$ від $x$.


Щоб побачити появу лінійної реакції орбіталей явно,
запишемо повну другу похідну енергії:

\begin{equation}
\frac{d^2E}{dx^2}
=
\frac{\partial^2 E}{\partial x^2}
+
2 \sum_i
\frac{\partial^2 E}{\partial x\,\partial p_i}\frac{dp_i}{dx}
+
\sum_{i,j}
\frac{\partial^2 E}{\partial p_i \partial p_j}
\frac{dp_i}{dx}\frac{dp_j}{dx}
+
\sum_i
\frac{\partial E}{\partial p_i}\frac{d^2 p_i}{dx^2}.
\label{eq:second_derivative_full}
\end{equation}



Останній член у \eqref{eq:second_derivative_full} зникає завдяки умові \eqref{eq:stationarity_p0}, і друга похідна набуває вигляду
\begin{equation}
\frac{d^2E}{dx^2}
=
\frac{\partial^2 E}{\partial x^2}
+
2 \sum_i
\frac{\partial^2 E}{\partial x\,\partial p_i}\frac{dp_i}{dx}
+
\sum_{i,j}
\frac{\partial^2 E}{\partial p_i \partial p_j}
\frac{dp_i}{dx}\frac{dp_j}{dx}.
\label{eq:second_derivative_simplified}
\end{equation}

У \emph{других} похідних з'являються складові $\partial p_i/\partial x$:

Оскільки рівність \eqref{eq:stationarity_p0} справедлива для всіх \(x\),
візьмемо повну похідну по \(x\) від лівої частини (тобто продиференціюємо по \(x\) рівняння стаціонарності):
\begin{equation}
\frac{d}{dx}\left[
\left.\frac{\partial E(x,\mathbf p)}{\partial p_i}\right|_{\mathbf p=\mathbf p^{(0)}(x)}
\right] = 0.
\label{eq:diff_stationarity_start}
\end{equation}

Розкриваємо повну похідну:
\begin{equation}
\frac{d}{dx}\left[\frac{\partial E}{\partial p_i}\right]
=
\left.\frac{\partial^2 E}{\partial x\,\partial p_i}\right|_{\mathbf p=\mathbf p^{(0)}(x)}
+
\sum_j
\left.\frac{\partial^2 E}{\partial p_i \partial p_j}\right|_{\mathbf p=\mathbf p^{(0)}(x)}
\frac{d p_j(x)}{dx}.
\label{eq:diff_stationarity_explicit}
\end{equation}

Оскільки ліва частина рівняння \eqref{eq:diff_stationarity_start} дорівнює нулю,
отримуємо
\begin{equation}
\left.\frac{\partial^2 E}{\partial p_i \partial p_j}\right|_{\mathbf p=\mathbf p^{(0)}(x)}
\frac{d p_j(x)}{dx}
= -
\left.\frac{\partial^2 E}{\partial x\,\partial p_i}\right|_{\mathbf p=\mathbf p^{(0)}(x)}.
\label{eq:intermediate_CPHF}
\end{equation}

Позначимо матрицю других похідних по параметрах (\emph{електронний гессіан} у просторі параметрів)
\begin{equation}
H_{ij}(x) \equiv
\left.\frac{\partial^2 E(x,\mathbf p)}{\partial p_i \partial p_j}\right|_{\mathbf p=\mathbf p^{(0)}(x)},
\label{eq:H_matrix_def}
\end{equation}
і вектор правої частини (змішана друга похідна)
\begin{equation}
g_i^{(x)} \equiv
\left.\frac{\partial^2 E(x,\mathbf p)}{\partial x\,\partial p_i}\right|_{\mathbf p=\mathbf p^{(0)}(x)}.
\label{eq:g_vector_def}
\end{equation}

Тоді рівняння \eqref{eq:intermediate_CPHF} компактно записується як
\begin{equation}
\sum_j H_{ij}(x)\;\frac{d p_j^{(0)}(x)}{dx} = -\,g_i^{(x)}.
\label{eq:CPHF_derived_compact}
\end{equation}

Якщо матриця \(H(x)\) невироджена (має обернену), то формально
\begin{equation}
\frac{d\mathbf p^{(0)}(x)}{dx} = -\,\mathbf{H}^{-1}(x)\,\mathbf g^{(x)}.
\label{eq:dpdx_solution}
\end{equation}

Отримане рівняння є системою, що має назву
зв'язано-збуреної (\emph{Coupled--Perturbed}) у загальному вигляді.
Воно визначає лінійну реакцію оптимальних варіаційних параметрів
на малу зміну зовнішнього параметра \(x\).

Конкретний вигляд цієї системи залежить від того, які саме параметри \(p_i\)
використано для задання орбіталей.
У методах \texttt{HF} та \texttt{DFT} параметри \(p_i\) відповідають
елементам матриці орбітальних обертань, а величини \(dp_i/dx\)
є їхньою лінійною реакцією.
У цьому випадку матриця \(H\) і вектор \(g^{(x)}\) набувають
звичних аналітичних форм, що приводить до стандартних рівнянь
\texttt{CPHF}/\texttt{CPKS}.


%% ========================================================
\section{Аналітичні похідні в методі Хартрі-Фока}
%% ========================================================

У попередньому розділі ми розглянули загальний варіаційний підхід:
енергію $E(x,\mathbf p)$, оптимальні параметри $\mathbf p^{(0)}(x)$,
електронний гессіан $H$ та рівняння \eqref{eq:dpdx_solution}
яке формально описує, як змінюються оптимальні варіаційні параметри
при малому збуренні зовнішнього параметра $x$.

\subsection{Орбітальні параметри в методі Хартрі--Фока}

У \texttt{HF} хвильова функція --- це детермінант Слейтера, повністю
визначений набором молекулярних орбіталей
\[
\phi_p = \sum_\mu C_{\mu p}\,\chi_\mu .
\]
Проте варіювати матрицю коефіцієнтів $\mathbf C$ безпосередньо неможливо через
вимогу ортонормованості орбіталей:
\[
\langle\phi_p|\phi_q\rangle=\delta_{pq}.
\]

Тому варіацію формулюють через \emph{нескінченно малі орбітальні обертання}. У детермінанті фізично значимими є лише ті обертання, що змішують \emph{зайняті} та \emph{віртуальні} орбіталі: обертання всередині підпросторів зайнятих або віртуальних МО не змінюють хвильову функцію, а лише переставляють базисні функції всередині однаково заселених підпросторів.


Допустима зміна зайнятої орбіталі має вигляд:
\begin{equation}\label{eq:small_rotation_intro}
\phi_i \;\longrightarrow\;
\phi_i^{(0)} + \sum_a U_{a i}\,\phi_a^{(0)},
\end{equation}
де змішуються лише зайняті $\phi_i$ з віртуальними $\phi_a$,
а матриця орбітальних обертань має антисиметричні блоки:
\[
U_{ai} = -U_{ia}.
\]

Така заміна означає ось що зайнята орбіталь під дією збурення трохи «втрачає свою форму» і набуває легку домішку віртуальних орбіталей. Віртуальні МО відіграють роль можливих «напрямків деформації», у які електронна хмара може відхилитися, реагуючи на зовнішній вплив. В малих базисах часто спостерігають занижену поляризовність, погану відтворюваність похідних енергії. Чим багатший базис, тим повніший набір можливих напрямків деформації й тим адекватніше описується фізична відповідь електронів.

\paragraph{Варіаційні параметри.}

У підході орбітальних обертань саме коефіцієнти
\begin{equation}\label{eq:p_HF_definition}
p_{ai} \equiv U_{ai}
\end{equation}
виступають справжніми варіаційними параметрами. Вони кількісно описують,
наскільки кожна зайнята орбіталь допускає домішку конкретної віртуальної,
тобто наскільки електронна хмара може деформуватися в цьому напрямку.
Таким чином, вектор параметрів
\[
\mathbf p = \{U_{ai}\}
\]
є компактним описом усіх фізично значущих змін хвильової функції
в межах \texttt{HF}-варіації.


\paragraph{Похідні параметрів.}
Коли зовнішній параметр $x$ змінюється на $dx$, оптимальні орбіталі
переходять у нові оптимальні орбіталі для $x+dx$.
Лінійна зміна МО має вигляд
\begin{equation}\label{eq:HF_MO_derivative}
\frac{d\phi_i}{dx}
=
\sum_a U_{ai}^{(x)} \phi_a ,
\end{equation}
що визначає коефіцієнти
\[
U_{ai}^{(x)} \equiv \frac{dU_{ai}}{dx}.
\]

Отже, у \texttt{HF} величини
\[
\frac{dp_{ai}}{dx}
=
\frac{dU_{ai}}{dx}
\equiv U_{ai}^{(x)}
\]
є \emph{лінійною реакцією зайнятих орбіталей на збурення $x$},
тобто тим, як зайняті МО змішуються з віртуальними при зміні зовнішнього параметра.

\paragraph{Рівняння CPHF.}
Матриця $\mathbf H$ у просторі параметрів $U_{ai}$ набуває вигляду
\begin{equation}\label{eq:H_ai_bj_def}
H_{ai,bj}
=
\left.\frac{\partial^2 E}{\partial U_{ai}\,\partial U_{bj}}\right|_{\mathbf U= 0}.
\end{equation}

Аналогічно, компоненти вектора змішаної похідної:
\begin{equation}\label{eq:g_ai_def}
g_{ai}^{(x)}
=
\frac{\partial }{\partial U_{ai}}\left.\frac{\partial E}{\partial x}\right|_{x = 0, \mathbf U= 0}.
\end{equation}

Після цих підстановок рівняння \eqref{eq:dpdx_solution} набуває вигляду
\begin{equation}\label{eq:CPHF_matrix_general}
\sum_{b j} H_{ai,bj}\;U_{bj}^{(x)} = -g_{ai}^{(x)}.
\end{equation}

Це і є загальна форма \emph{зв'язано--збурених рівнянь} (\texttt{CPHF}).
Вони визначають реакцію МО на довільне збурення $x$.

\begin{commentbox}[Приклад]
Розглянемо спрощену систему, де є одна зайнята орбіталь $i$ та одна
віртуальна орбіталь $a$, розкладені за тими самими базисними функціями
$\{\chi_\mu\}$:
\[
\phi_i^{(0)} = C_{1i}\chi_1 + C_{2i}\chi_2 + C_{3i}\chi_3,\qquad
\phi_a^{(0)} = C_{1a}\chi_1 + C_{2a}\chi_2 + C_{3a}\chi_3.
\]

Після нескінченно малого обертання орбіталі змінюються за правилом
\eqref{eq:small_rotation_intro}:
\[
\phi_i = \phi_i^{(0)} + U_{ai}\phi_a^{(0)}.
\]

Якщо система залежить від зовнішнього параметра $x$, то в першому порядку відповідь орбіталі записується як
\[
\phi_i(x) \approx \phi_i^{(0)} + x\sum_a U_{ai}^{(x)}\,\phi_a^{(0)},
\]
де $U_{ai}^{(x)}\equiv dU_{ai}/dx$ — саме ті похідні, що даються рівняннями CPHF.

Підставивши явні розклади у базисі $\{\chi_\mu\}$, отримуємо рівняння для коефіцієнтів:
\[
C_{\mu i}(x) \approx C_{\mu i} + x\sum_a U_{ai}^{(x)}\,C_{\mu a},
\qquad \mu=1,2,3.
\]

Отже, похідні $U_{ai}^{(x)}$ описують, як при малому змінені $x$ змінюються
коефіцієнти зайнятої орбіталі в напрямку коефіцієнтів віртуальної орбіталі.
\end{commentbox}



\subsection{Стандартна форма рівнянь CPHF}

У методі \texttt{HF} електронний гессіан у просторі $U_{ai}$ має вигляд
\begin{equation}\label{eq:H_standard}
H_{ai,bj}
=
(\varepsilon_a - \varepsilon_i)\,\delta_{ab}\delta_{ij}
+ (ai||bj),
\end{equation}
де $(ai||bj)$ --- антисиметризовані двоелектронні інтеграли.

Праві частини $g_{ai}^{(x)}$ залежать від збурення $x$.
Наприклад, при ядерному русі це:
\[
g_{ai}^{(x)} = \langle \phi_a | \hat F^{(x)} | \phi_i \rangle,
\]
де $\hat h^{(x)}$ --- похідна одноелектронного гамільтоніана,
а $\hat V^{(x)}_{ee}$ --- похідна фоківського потенціалу.

Таким чином, рівняння \eqref{eq:CPHF_matrix_general} у явному вигляді стає
\begin{equation}\label{eq:CPHF_explicit}
(\varepsilon_a - \varepsilon_i)\,U_{ai}^{(x)}
+ \sum_{b j} (ai||bj)\,U_{bj}^{(x)}
= -\,g_{ai}^{(x)}.
\end{equation}

Це --- стандартні рівняння \texttt{CPHF}/\texttt{CPKS}, які визначають
лінійну реакцію молекулярних орбіталей при довільному збуренні $x$.


Розв'язавши \eqref{eq:CPHF_explicit}, отримуємо $U_{ai}^{(x)}$,
тобто як зайняті орбіталі <<змішуються>> з віртуальними у відповідь на збурення.

\begin{commentbox}
Розв'язок рівнянь CPHF (Coupled Perturbed Hartree-Fock) зазвичай є чисельним. У практичних обчисленнях квантової хімії ці рівняння представляють собою велику систему лінійних рівнянь виду $\mathbf{H} \mathbf{U}^{(x)} = -\mathbf{g}^{(x)}$, де матриця $\mathbf{H}$ (електронний гессіан) може мати високу розмірність (наприклад, тисячі або мільйони елементів для реальних молекул).

\begin{itemize}
\item Прямі методи (наприклад, інвертування матриці) використовуються рідко, тільки для малих систем, оскільки вони обчислювально дорогі ($O(N^3)$ складність, де $N$ --- розмірність).
\item Ітераційні методи є стандартними для великих систем: наприклад, метод спряжених градієнтів (conjugate gradient), DIIS (Direct Inversion in the Iterative Subspace) або інші солвери для лінійних рівнянь. Вони дозволяють уникнути повного зберігання матриці $\mathbf{H}$ в пам'яті, обчислюючи її добутки з векторами <<на льоту>>.
\end{itemize}

Аналітичний розв'язок можливий лише для тривіальних або модельних систем (наприклад, атом водню), але в реальних розрахунках все робиться чисельно з контролем збіжності ітерацій. Якщо збурень багато.
\end{commentbox}

\begin{figure}[h!]
\centering
\begin{tikzpicture}[
    scale=2,
    every node/.style={font=\scriptsize, very thin},
    node distance=1cm,
    startstop/.style={rectangle, rounded corners,
        minimum width=6cm, minimum height=0.6cm,
        text centered, draw=black, fill=red!30},
    process/.style={rectangle,
        minimum width=6cm, minimum height=0.6cm,
        text centered, draw=black, fill=orange!30},
    io/.style={trapezium, trapezium left angle=70,
        trapezium right angle=110, minimum height=0.6cm,
        text centered, draw=black, fill=blue!30,
        minimum width=2cm},
    decision/.style={diamond, aspect=2,
        text centered, draw=black, fill=green!30},
]

% Блоки
\node (start) [io]
{Початкове наближення $U_{ai}^{(x)}=0$};

\node (buildRHS) [process, below=0.6cm of start, text width=5cm]
{Обчислення правої частини: \\[2pt]
$g_{ai}^{(x)} = \displaystyle
\left.\frac{\partial^2E}{\partial x\,\partial U_{ai}}\right|_0 $
\\[2pt]
(похідна Фока по $x$)};

\node (buildH) [process, below=0.6cm of buildRHS, text width=5cm]
{Побудова електронного гессіана:\\[2pt]
$H_{ai,bj} = (\varepsilon_a-\varepsilon_i)\delta_{ab}\delta_{ij}
+ (ai||bj)$};

\node (solveU) [process, below=0.6cm of buildH, text width=5cm]
{Розв’язання матричних рівнянь CPHF:\\[4pt]
$\displaystyle
\sum_{bj} H_{ai,bj}\,U_{bj}^{(x)}
= -g_{ai}^{(x)}$};

\node (buildD1) [process, below=0.6cm of solveU, text width=5cm]
{Оновлення похідної густини:\\[4pt]
$\displaystyle
\mathbf D^{(1)}_{\!ov} = U^{(x)}, \qquad
\mathbf D^{(1)}_{\!vo} = U^{(x)\dagger}$};

\node (check) [decision, below=0.6cm of buildD1]
{$\|U^{(x)}_{n+1}-U^{(x)}_{n}\|\le\delta$};

\node (props) [startstop, below=0.8cm of check, text width=10cm]
{Використання $\mathbf D^{(1)}$ або $U^{(x)}$ \\[2pt]
для обчислення другої похідної енергії, \\[2pt]
тензорів поляризовності, Раман-поляризовності та гессіанів.};

% Стрілки
\draw[-latex] (start) -- (buildRHS);
\draw[-latex] (buildRHS) -- (buildH);
\draw[-latex] (buildH) -- (solveU);
\draw[-latex] (solveU) -- (buildD1);
\draw[-latex] (buildD1) -- (check);

% Петля (Ні)
\draw[-latex] (check.west) -- node[above]{Ні} +(-2,0)
    coordinate (loop) -- (loop |- solveU.west)
    |- (solveU.west);

% Так — вихід
\draw[-latex] (check.east) -- node[above]{Так} +(2,0)
    coordinate (out) |- (props.east);

\end{tikzpicture}
\caption{Процедура розв`язку рівнянь \texttt{CPHF}.
}
\end{figure}


%\section{Похідні інтегралів і позначення}
%
%Позначення похідних одно- і двоелектронних інтегралів по параметру \(x\):
%\begin{align}
%	h_{\mu\nu}^{(x)}             & \equiv \frac{\partial h_{\mu\nu}}{\partial x},
%	\qquad h_{\mu\nu}^{(xy)} \equiv \frac{\partial^2 h_{\mu\nu}}{\partial x\partial y},
%	\label{eq:h_deriv_def}                                                                    \\
%	(\mu\nu|\lambda\sigma)^{(x)} & \equiv \frac{\partial (\mu\nu|\lambda\sigma)}{\partial x},
%	\qquad \text{і т. д.}
%	\label{eq:eri_deriv_def}
%\end{align}
%
%Матриця перекриття та її похідні:
%\begin{equation}
%	S_{\mu\nu}^{(x)} \equiv \frac{\partial S_{\mu\nu}}{\partial x},\qquad
%	S_{\mu\nu}^{(xy)} \equiv \frac{\partial^2 S_{\mu\nu}}{\partial x\partial y}.
%	\label{eq:S_deriv_def}
%\end{equation}
%
%MO-індекси: \(i,j\) — зайняті, \(a,b\) — віртуальні; \(C_{\mu p}\) — коефіцієнти MO.

%% ========================================================
\subsection{Градієнт (перша похідна) — явний вираз у методі HF}
%% ========================================================

У методі \texttt{HF} умова стаціонарності \eqref{eq:stationarity_p0} набуває вигляду
\[
\left.\frac{\partial E}{\partial U_{ai}}\right|_{U=U^{(0)}} = 0.
\]

Отже, перша похідна енергії визначається \emph{лише} явною залежністю
гамільтоніана від $x$:
\begin{equation}\label{eq:HF_gradient_basic}
\frac{dE}{dx}
=
\left.\frac{\partial E}{\partial x}\right|_{U=U^{(0)}} .
\end{equation}

\paragraph{Одноелектронний внесок.}
У \texttt{HF} енергія має вигляд
\[
E = \sum_i h_{ii} + \tfrac{1}{2}\sum_{ij}\big( J_{ij} - K_{ij} \big),
\]
де $h_{pq}$ --- одноелектронна частина, $J_{ij}$ і $K_{ij}$ --- кулонівський
та обмінний інтеграли.

Якщо зовнішній параметр $x$ змінює ядерний потенціал або інший зовнішній
оператор, то перша похідна енергії включає зміну тільки відповідних
інтегралів:
\begin{equation}\label{eq:HF_gradient_one_electron}
\frac{dE}{dx}
=
\sum_i \frac{\partial h_{ii}}{\partial x}
+
\tfrac12 \sum_{ij}
\left(
\frac{\partial J_{ij}}{\partial x}
-
\frac{\partial K_{ij}}{\partial x}
\right),
\end{equation}
оскільки зміна коефіцієнтів МО (тобто зміна $U_{ai}$) не дає внеску
в першу похідну через стаціонарність.

\paragraph{Форма через матрицю густини.}
У більш компактній матричній формі:
\begin{equation}\label{eq:HF_gradient_density_matrix}
\frac{dE}{dx}
=
\mathrm{Tr}\!\left[\,\mathbf D\,\frac{\partial \mathbf h}{\partial x}\right]
+
\tfrac12\,\mathrm{Tr}\!\left[\mathbf D\,\frac{\partial (\mathbf J - \mathbf K)}{\partial x}\right],
\end{equation}
де $\mathbf D$ --- матриця густини у базисі АО.

Це класичний вираз градієнта у \texttt{HF}, який використовується для
обчислення сил, геометричних градієнтів, похідних по зовнішньому полі тощо.

\paragraph{Важлива властивість.}
Завдяки умові стаціонарності, у градієнті \emph{не з'являються} похідні
орбітальних обертань $dU_{ai}/dx$.
Інакше кажучи, у \texttt{HF} перша похідна енергії не потребує розв’язку
рівнянь лінійної реакції. Всі складні терміни, пов'язані з реакцією МО,
виникають лише у другій похідній (у наступному розділі), де вони приводять
до системи \texttt{CPHF}.

% --------------------------

%Перша повна похідна енергії по параметру \(x\) у HF має вигляд
%(усі похідні інтегралів беруться при фіксованих орбіталях; вклад від зміни орбіталей зникає
%через варіаційну умову):
%\begin{equation}
%	\frac{dE_{\mathrm{HF}}}{dx}
%	=
%	\mathrm{Tr}\!\big[\mathbf D\,\mathbf h^{(x)}\big]
%	+ \tfrac12 \mathrm{Tr}\!\big[\mathbf D\,\mathbf G^{(x)}\big]
%	- \mathrm{Tr}\!\big[\mathbf F\,\mathbf D\,\mathbf S^{(x)}\big]
%	+ \frac{dV_{\mathrm{NN}}}{dx}.
%	\label{eq:HF_gradient}
%\end{equation}
%
%Тут \(\mathbf F = \mathbf h + \mathbf G\) — матриця Фока, \(\mathbf D\) — матриця густини у AO-базисі,
%а \(\mathbf G^{(x)}\) — похідна двоелектронного внеску (через похідні ERI).
%
%Рівняння \eqref{eq:HF_gradient} — те, що реалізовано у кодах як аналітичний градієнт.

%% ========================================================
\subsection{Друга похідна — явний вигляд у методі HF}
%% ========================================================

У попередніх підрозділах ми показали загальний варіаційний результат
\[
\frac{d^2E}{dx^2}
=
\frac{\partial^2 E}{\partial x^2}
+2\sum_{ai}\frac{\partial^2 E}{\partial x\,\partial U_{ai}}\frac{dU_{ai}}{dx}
+\sum_{ai}\sum_{bj}\frac{\partial^2 E}{\partial U_{ai}\partial U_{bj}}
\frac{dU_{ai}}{dx}\frac{dU_{bj}}{dx}.
\]
Тепер випишемо цю формулу в термінах вже введених об'єктів
$g^{(x)}_{ai}$ та $H_{ai,bj}$ і одержимо робочу форму для другої похідної.


Остаточний вираз для другої похідної набуває компактної форми
\begin{equation}\label{eq:d2E_final_boxed}
\displaystyle
\frac{d^2E}{dx^2}
=
\left.\frac{\partial^2 E}{\partial x^2}\right|_{U}
-
\sum_{ai}\sum_{bj}
g^{(x)}_{ai}\,(\mathbf{H}^{-1})_{ai,bj}\,g^{(x)}_{bj}.
\end{equation}

Цю формулу часто записують у скороченому матричному вигляді:
\[
\frac{d^2E}{dx^2} = \left.\frac{\partial^2 E}{\partial x^2}\right|_{U} - \mathbf g^{(x)\,\mathrm T} \mathbf{H}^{-1}\mathbf g^{(x)}.
\]

\paragraph{Інтерпретація і розгорнуті компоненти.}
\begin{itemize}
    \item Перший член $\displaystyle\left.\frac{\partial^2 E}{\partial x^2}\right|_{U}$ — пряма (явна) друга похідна енергії по $x$ при фіксованих орбіталях. У матричній формі її можна виписати як слід:
\begin{multline*}
    \left.\frac{\partial^2 E}{\partial x^2}\right|_{U}
    =
    \mathrm{Tr}\!\left[\mathbf D\,\frac{\partial^2\mathbf h}{\partial x^2}\right]
    +\tfrac12\,\mathrm{Tr}\!\left[\mathbf D\,\frac{\partial^2(\mathbf J-\mathbf K)}{\partial x^2}\right]
    + \\ +\text{додаткові Pulay-терміни явного другого порядку}.
\end{multline*}
    \item Другий член $-\mathbf g^{\mathrm T} H^{-1}\mathbf g$ — вклад, що походить від лінійної реакції орбіталей. Він завжди від'ємний для позитивно визначного $\mathbf H$ і фізично відповідає зниженню енергії за рахунок релаксації орбіталей під дією збурення.
    \item Компоненти $g^{(x)}_{ai}$ у випадку \texttt{HF} зазвичай мають вигляд
    \[
    g^{(x)}_{ai} = \langle \phi_a | \hat h^{(x)} | \phi_i \rangle
    + \langle \phi_a | \hat V^{(x)}_\text{HF} | \phi_i \rangle,
    \]
    де $\hat h^{(x)}$ — явна похідна одноелектронного оператора (наприклад, похідна ядерного потенціалу), а $\hat V^{(x)}_\text{HF}$ — похідна фоківського потенціалу (вона містить похідні двоелектронних інтегралів).
\end{itemize}

\paragraph{Примітки щодо практики обчислення.}
\begin{enumerate}
    \item Якщо базис залежить від параметра $x$ (наприклад, при зміщенні ядер у атомно-орбітальному базисі), то необхідно врахувати Pulay-терміни як у першій, так і у другій похідній; вони входять у явну частину $\partial^2E/\partial x^2$ і в компоненту $g^{(x)}$ (через похідні інтегралів).
    \item У випадку виродженості або слабко визначеного H потрібно проектувати фізичний підпростір обертань (ці обернення, що фізично змінюють хвильову функцію) і інвертувати H в цьому підпросторі.
    \item Формула \eqref{eq:d2E_final_boxed} є базовою для обчислення других похідних: матриці силових констант, поляризованості 2-го порядку, динамічних поправок тощо.
\end{enumerate}


%% ========================================================
\subsection{Змішана похідна у методі HF — явний вигляд}
%% ========================================================

Розглянемо два зовнішні параметри $x$ і $y$. Для методу \texttt{HF} змішана похідна
\[
\frac{d^2E}{dx\,dy}
\]
має ту саму структурну форму, що й у загальному варіаційному випадку, але всі складові можна виписати явно.

\paragraph{1. Явна (пряма) часткова друга похідна при фіксованих орбіталях.}
При фіксованих орбіталях ($U$ константні)
\begin{multline}\label{eq:explicit_partial_xy}
\left.\frac{\partial^2 E}{\partial x\,\partial y}\right|_{U}
=
\mathrm{Tr}\!\left[\mathbf D\,\frac{\partial^2\mathbf h}{\partial x\,\partial y}\right]
+\tfrac12\,\mathrm{Tr}\!\left[\mathbf D\,\frac{\partial^2(\mathbf J-\mathbf K)}{\partial x\,\partial y}\right]
+ \\ + \text{(Pulay-доданки при залежному базисі)}.
\end{multline}
Тут $\mathbf D$ — матриця густини у АО-базисі, $\mathbf h$ — матриця одноелектронного оператора, $\mathbf J,\mathbf K$ — відповідні матриці кулон- та обмінних операторів.

\paragraph{2. Вектори змішаних похідних \(g^{(x)}\) та \(g^{(y)}\).}
Для \texttt{HF} компоненти вектора $g^{(x)}$ мають класичний вигляд
\begin{equation}\label{eq:g_ai_HF_explicit}
g^{(x)}_{ai}
=
\langle \phi_a \,|\, \tfrac{\partial\hat h}{\partial x} \,|\, \phi_i \rangle
+
\langle \phi_a \,|\, \tfrac{\partial\hat V_{\mathrm{HF}}}{\partial x} \,|\, \phi_i \rangle
+
\text{(Pulay-доданку від залежного базису)},
\end{equation}
де
\[
\hat V_{\mathrm{HF}} = \hat J - \hat K
\]
— фоківський оператор. Аналогічно для $g^{(y)}_{ai}$:
\[
g^{(y)}_{ai}
=
\langle \phi_a \,|\, \tfrac{\partial\hat F}{\partial y} \,|\, \phi_i \rangle
+ \text{(Pulay)}.
\]

Якщо базис залежить від параметрів ($\chi_\mu=\chi_\mu(x,y,\dots)$), то в $g^{(x)}_{ai}$ додатково входять члени, пропорційні похідним перекриттів $\mathbf S^{(x)}$ і похідним інтегралів; ці Pulay-доданки можна виписати стандартно через $\mathbf F$ і $\mathbf S^{(x)}$.



\paragraph{3. Рівняння лінійної реакції і розв'язок.}
Зв'язані рівняння для реакції при збуренні $y$:
\begin{equation}\label{eq:CPHF_HF_xy}
\sum_{bj} H_{ai,bj}\,\frac{dU_{bj}}{dy} = -\,g^{(y)}_{ai},
\end{equation}
і при невиродженому $H$ розв'язок
\[
\frac{dU_{ai}}{dy} = -\sum_{bj}(H^{-1})_{ai,bj}\,g^{(y)}_{bj}.
\]

\paragraph{5. Остаточна формула для змішаної похідної у HF.}
Підставивши розв'язок у загальну формулу \eqref{eq:mixed_start}, одержуємо
\begin{equation}\label{eq:mixed_HF_final}
\frac{d^2E}{dx\,dy}
=
\mathrm{Tr}\!\left[\mathbf D\,\frac{\partial^2\mathbf h}{\partial x\,\partial y}\right]
+\tfrac12\,\mathrm{Tr}\!\left[\mathbf D\,\frac{\partial^2(\mathbf J-\mathbf K)}{\partial x\,\partial y}\right]
-
\sum_{ai}\sum_{bj} g^{(x)}_{ai}\,(\mathbf{H}^{-1})_{ai,bj}\,g^{(y)}_{bj}
\end{equation}
або компактно
\[
\frac{d^2E}{dx\,dy}
=
\left.\frac{\partial^2 E}{\partial x\,\partial y}\right|_{U}
-
\mathbf g^{(x)\,\mathrm T} H^{-1}\mathbf g^{(y)}.
\]

\paragraph{6. Коментарі щодо Pulay-термінів і практики.}
\begin{itemize}
  \item Якщо базис залежить від $x$ або $y$, то перший член у формулі \eqref{eq:mixed_HF_final} повинен включати Pulay-члени першого та другого порядку. Зокрема, у частковій другій похідній при фіксованих орбіталях входять члени з $\partial^2\mathbf h/\partial x\partial y$ та члени виду $\mathrm{Tr}[\mathbf F\,\partial^2\mathbf S/\partial x\partial y]$ (або еквівалентні вирази).
  \item Вектори $\mathbf g^{(x)}$ і $\mathbf g^{(y)}$ також містять похідні двоелектронних інтегралів — їх треба збирати при побудові правих частин CPHF.
  \item Для обчислення кількох змішаних похідних доцільно однаково розкладати або факторизувати матрицю $H$ (наприклад, LU/Cholesky/диагоналізація в робочому підпросторі) і потім для кожного збурення швидко знаходити $H^{-1}\mathbf g^{(z)}$ шляхом прямого чи зворотного ходу.
  \item Симетрія змішаної похідної ($d^2E/dxdy = d^2E/dydx$) гарантується при правильному урахуванні всіх Pulay-термінів і при точному розв'язанні CPHF.
\end{itemize}


%% ============================================================
\section{Поляризовність у методі HF}
%% ============================================================

Поляризовність визначається як друга похідна енергії за компонентами
зовнішнього електричного поля:
\begin{equation}\label{eq:alpha_def}
\alpha_{\mu\nu}
=
-\left.\frac{\partial^2 E}{\partial \mathcal E_\mu\,\partial \mathcal E_\nu}\right|_{\mathcal E = 0}.
\end{equation}

У присутності однорідного електричного поля одноелектронний оператор
набуває вигляду
\[
\hat h(\boldsymbol{\mathcal E})
=
\hat h_0 - \sum_{\mu} \mathcal E_\mu\,\hat\mu_\mu,
\]
тому
\[
\frac{\partial \hat h}{\partial \mathcal E_\mu} = -\hat\mu_\mu,
\qquad
\frac{\partial^2 \hat h}{\partial \mathcal E_\mu\,\partial \mathcal E_\nu}=0.
\]

% ------------------------------------------------------------

Похідна фокіана по компоненті поля \(\mu\) збігається з похідною одноелектронної частини:
\[
\hat F^{(\mu)} = \frac{\partial \hat F}{\partial \mathcal E_\mu}
= -\hat\mu_\mu.
\]

Звідси елементи вектора збурення у рівняннях \texttt{CPHF}:
\begin{equation}\label{eq:HF_g_vector_field}
g^{(\mu)}_{ai}
=
\left.\frac{\partial^2 E}{\partial \mathcal E_\mu\,\partial U_{ai}}\right|_{\mathcal E=0}
=
F^{(\mu)}_{ai}
=
-(\mu_{\mu})_{ai},
\end{equation}
де \(\mu = x,y,z\) - компонента поля (або дипольного моменту) вздовж відповідної осі.

Розв’язавши рівняння \texttt{CPHF},
\[
\sum_{bj} H_{ai,bj}\,U^{(\mu)}_{bj} = -g^{(\mu)}_{ai},
\]
отримуємо орбітальну відповідь \(U^{(\mu)}_{ai}\).

% ------------------------------------------------------------

Перший порядок густини:
\[
D^{(\mu)}_{ai} = 2\,U^{(\mu)}_{ai},
\qquad
D^{(\mu)}_{ia} = 2\,U^{(\mu)}_{ai},
\qquad
D^{(\mu)}_{ij}=0,\quad
D^{(\mu)}_{ab}=0.
\]

Оскільки похідна Фока за компонентами поля лінійна й дорівнює
\[
F^{(\nu)}_{pq} = -(\mu_{\nu})_{pq},
\]
змішана друга похідна енергії у RHF набуває вигляду:
\begin{equation}\label{eq:HF_mixed_second_derivative}
\frac{\partial^2 E}{\partial \mathcal E_\mu\,\partial \mathcal E_\nu}
=
\mathrm{Tr}\!\left[D^{(\mu)} F^{(\nu)}\right]
=
-4 \sum_{ai}
U^{(\mu)}_{ai}\,(\mu_{\nu})_{ai}.
\end{equation}

Підставляючи у означення поляризовності~\eqref{eq:alpha_def},
маємо остаточний симетричний результат:
\begin{equation}\label{eq:HF_polarizability_final_sym}
\alpha_{\mu\nu}
=
4 \sum_{ai}
U^{(\mu)}_{ai}\,(\mu_{\nu})_{ai}
\qquad
\alpha_{\mu\nu}=\alpha_{\nu\mu}.
\end{equation}

Поляризовність у HF --- це подвійний проєкційний добуток між лінійною
орбітальною реакцією системи на поле \(\mathcal E_\mu\)
та матрицею дипольного моменту в напрямку \(\nu\).


\subsection{Детальний розбір повної CPHF-прцедури для знаходження тензора поляризовності}

В файлі \inputcode{h2_custom_CPHF.py} подано покроковий структурований опис коду, який реалізує обчислення тензора
поляризовності молекули \ce{H2O} за допомогою \texttt{CPHF} на базі самостійної побудови
електронного гессіана та правих частин рівнянь.

% ============================================================
\paragraph{1. SCF, МО та базові структурні параметри.}

Початковий етап: будуємо молекулу, виконуємо \texttt{RHF} та отримуємо
\(
 \mathbf C,\;
 \varepsilon_i,\;
 n_{\text{occ}},\;
 n_{\text{vir}}
\),
де
\[
\mathbf F \mathbf C = \mathbf S\,\mathbf C\,\boldsymbol\varepsilon ,\qquad
\mathbf C=\{C_{pi}\}.
\]

\begin{minted}{python}
import numpy as np
from pyscf import gto, scf
from pyscf.prop.polarizability.rhf import Polarizability

# === H2O ===
mol = gto.M(
    atom='''
    O 0.000000 0.000000 0.000000
    H 0.000000 -0.757000 0.587000
    H 0.000000 0.757000 0.587000
    ''',
    basis='6-31g',
    verbose=0
)

mf = scf.RHF(mol).run(conv_tol=1e-12)
print(f"Енергія: {mf.e_tot:.10f} Hartree")

mo_coeff = mf.mo_coeff      # C_{p i}
mo_energy = mf.mo_energy    # енергії МО
mo_occ = mf.mo_occ          # заповнення МО

nocc = int(np.sum(mo_occ > 0))  # кількість зайнятих МО
nvir = len(mo_energy) - nocc    # кількість віртуальних МО
ov = nocc * nvir                # розмірність простору ia
\end{minted}

% ============================================================
\paragraph{2. Дипольні інтеграли в АО та перехід у МО.}

Будуються матриці
\( (\mu_x)_{pq}, (\mu_y)_{pq}, (\mu_z)_{pq}\)
та проєктуються на МО:
\[
(\mu_{\nu})_{ij} = C_{pi}\cdot (\mu_\nu)_{pq}\cdot C_{qj}.
\]

\begin{minted}{python}
with mol.with_common_orig((0,0,0)):
    dip_ao = mol.intor_symmetric('int1e_r', comp=3)  # (3, nao, nao)

dip_mo = np.einsum('pi,xpq,qj->xij', mo_coeff, dip_ao, mo_coeff)
\end{minted}

% ============================================================
\paragraph{3. Підготовка підпросторів $C_{\text{occ}}$, $C_{\text{vir}}$.}

\[
C_{\text{occ}} = C[:,\, i],\qquad
C_{\text{vir}} = C[:,\, a].
\]

\begin{minted}{python}
C_occ = mo_coeff[:, :nocc]
C_vir = mo_coeff[:, nocc:]
\end{minted}

% ============================================================
\paragraph{4. Двоелектронні інтеграли: AO → MO.}

Будуємо блоки
\((ia|jb)\), \((ij|ab)\), \((ib|ja)\),
які входять у CPHF-гессіан.

\begin{align*}
(ia|jb) &= C_{pi}^{\text{occ}} C_{qa}^{\text{vir}}
          C_{rj}^{\text{occ}} C_{sb}^{\text{vir}}\, (pq|rs),
\\
(ij|ab) &= C_{pi}^{\text{occ}} C_{qj}^{\text{occ}}
          C_{ra}^{\text{vir}} C_{sb}^{\text{vir}}\, (pq|rs).
\end{align*}

\begin{minted}{python}
eri_ao = mol.intor('int2e')

print("ERI → MO...")

# (ia|jb)
eri_iajb = np.einsum('pi,pqrs->iqrs', C_occ, eri_ao)
eri_iajb = np.einsum('qa,iqrs->iars', C_vir, eri_iajb)
eri_iajb = np.einsum('rj,iars->iajs', C_occ, eri_iajb)
eri_iajb = np.einsum('sb,iajs->iajb', C_vir, eri_iajb)

# (ij|ab)
eri_ijab = np.einsum('pi,pqrs->iqrs', C_occ, eri_ao)
eri_ijab = np.einsum('qj,iqrs->ijrs', C_occ, eri_ijab)
eri_ijab = np.einsum('ra,ijrs->ijas', C_vir, eri_ijab)
eri_ijab = np.einsum('sb,ijas->ijab', C_vir, eri_ijab)

# (ib|ja) = перестановка (ia|jb)
eri_ibja = eri_iajb.transpose(0, 3, 2, 1)
\end{minted}

% ============================================================
\paragraph{5. Побудова гессіана CPHF.}

Електронний гессіан:
\[
H_{ia,jb}
= (\varepsilon_a - \varepsilon_i)\delta_{ij}\delta_{ab}
+ 4(ia|jb) - (ij|ab) - (ib|ja).
\]

\begin{minted}{python}
print("H матриця...")

i_occ = np.arange(nocc).repeat(nvir)
a_vir = np.tile(np.arange(nvir), nocc)

H = np.diag(mo_energy[nocc + a_vir] - mo_energy[i_occ])

H += 4.0 * eri_iajb.reshape(ov, ov)
H -= eri_ijab.transpose(0,2,1,3).reshape(ov, ov)
H -= eri_ibja.transpose(0,3,2,1).reshape(ov, ov)

print(f"H: {H.shape}")
\end{minted}

% ============================================================
\paragraph{6. Права частина CPHF: блок дипольних матриць.}

\[
g^{(x)}_{ia} = - (\mu_x)_{ia} =
- \langle \phi_i \mid \hat{\mu}_x \mid \phi_a \rangle,.
\]

\begin{minted}{python}
g_ov = -dip_mo[:, :nocc, nocc:]
\end{minted}

% ============================================================
\paragraph{7. Розв'язання системи CPHF:}
\[
\mathbf{H}\,\mathbf{U}^{(x)} = -\mathbf{g}^{(x)}.
\]

\begin{minted}{python}
print("\nCPHF...")

U = np.zeros((3, nocc, nvir))
for x in range(3):
    U[x] = np.linalg.solve(H, -g_ov[x].ravel()).reshape(nocc, nvir)
\end{minted}

% ============================================================
\paragraph{8. Тензор поляризовності.}

Основна формула:
\[
\alpha_{\nu\mu}
=  4 \sum_{ai}
U^{(\mu)}_{ai}\, (\mu_\nu)_{ai}
\qquad
\alpha_{\nu\mu} = \alpha_{\mu\nu}.
\]

\begin{minted}{python}
mu_ov = dip_mo[:, :nocc, nocc:]

alpha = 4 * np.einsum('xia,yia->xy', mu_ov, U)
\end{minted}

\begin{commentbox}
\begin{itemize}
   \item  Тут використано пряму лінійну алгебру (\inlinecode{np.linalg.solve}) --- тобто жодної ітераційної процедури в коді немає: матриця \(H\) розв'язується одразу (алгоритмічно --- LU-розклад або гасіння, як реалізовано в NumPy).
  \item Матриця \(H\) має розмір \((n_\text{occ} n_\text{vir})^2\). Для великих молекул це швидко стає непрактичним; у реальних імплементаціях використовують розріджені або ітеративні методи, а також симетрії.
  \item Перетворення ERI в МО-базис є дорогим по пам'яті; часто будують лише блоки, як це зроблено вище, або використовують часткові трансформації (поступово).
\end{itemize}
\end{commentbox}



%% ------------------------------------------------------------
%\subsection{Структурована схема обчислення}
%% ------------------------------------------------------------
%
%\begin{enumerate}
%    \item Провести незбурений HF-розрахунок і отримати МО, орбітальні енергії
%          \(\varepsilon_p\), двоелектронні інтеграли \((pq|rs)\) та матриці
%          дипольного моменту \(\mu^{(\lambda)}\).
%
%    \item Для кожної компоненти поля \(\mu=x,y,z\) побудувати праву частину
%    \[
%      g^{(\mu)}_{ai} = -\mu^{(\mu)}_{ai}.
%    \]
%
%    \item Розв’язати рівняння \texttt{CPHF}
%    \[
%      \sum_{bj} H_{ai,bj}\,U^{(\mu)}_{bj} = -g^{(\mu)}_{ai},
%    \]
%    отримавши \(U^{(\mu)}_{ai}\).
%
%    \item Побудувати перший порядок густини
%    \[
%      D^{(\mu)}_{ai} = D^{(\mu)}_{ia} = U^{(\mu)}_{ai}.
%    \]
%
%    \item Обчислити змішані другі похідні енергії
%    \[
%      \frac{\partial^2 E}{\partial \mathcal E_\mu\,\partial \mathcal E_\nu}
%      =
%      -2\sum_{ai} U^{(\mu)}_{ai}\,\mu^{(\nu)}_{ai}.
%    \]
%
%    \item Поляризовність
%    \[
%      \alpha_{\mu\nu}
%      =
%      2\sum_{ai} U^{(\mu)}_{ai}\,\mu^{(\nu)}_{ai}.
%    \]
%\end{enumerate}


%% ============================================================
\section{Ядерний гесіан у методі HF}
%% ============================================================

Ядерний гесіан --- це матриця других похідних повної електронної
енергії за двома довільними ядерними координатами $x=\vect{R}_A$ і $y=\vect{R}_B$:
\begin{equation}\label{eq:HF_Hessian_def}
H_{xy}
=
\frac{d^2 E}{dx\,dy}.
\end{equation}

У методі Гартрі--Фока енергія залежить від координат ядер двома способами:
(1) явно — через інтеграли одно- і двоелектронної частин гамільтоніана;
(2) неявно — через матрицю густини, яка змінюється при збуренні геометрії.
Тому друга похідна завжди розкладається на явну і неявну частини:
\begin{equation}\label{eq:HF_Hessian_split}
\frac{d^2E}{dx\,dy}
=
\left.
\frac{\partial^2E}{\partial x\,\partial y}
\right|_{\mathbf D}
+
\mathrm{Tr}
\!\left[
\frac{d\mathbf D}{dx}\,
\frac{\partial \mathbf F}{\partial y}
\right].
\end{equation}

% ------------------------------------------------------------
\paragraph{Явна частина.}
% ------------------------------------------------------------

Явна частина містить похідні інтегралів за ядерними координатами:
\begin{equation}
\left.
\frac{\partial^2E}{\partial x\,\partial y}
\right|_{\mathbf D}
=
\mathrm{Tr}\!\left[
\mathbf D\,\frac{\partial^2 \mathbf h}{\partial x\,\partial y}
\right]
+
\frac12
\mathrm{Tr}\!\left[
\mathbf D\,\frac{\partial^2 (\mathbf J - \mathbf K)}{\partial x\,\partial y}
\right].
\end{equation}
Це чисто інтегральні (геометричні) внески.

% ------------------------------------------------------------
\paragraph{Неявна частина і роль CPHF.}
% ------------------------------------------------------------

Неявна частина визначається реакцією МО на геометричне збурення.
Похідна густини задається амплітудами обертань $U_{ai}^{(x)}$, отриманими
з рівнянь \texttt{CPHF}:
\begin{equation}\label{eq:HF_Hessian_CPHF}
\sum_{bj}
H_{ai,bj}\,U^{(x)}_{bj}
=
-g^{(x)}_{ai}.
\end{equation}

Перший порядок густини:
\begin{equation}
D^{(x)}_{ai}=U^{(x)}_{ai},\qquad
D^{(x)}_{ia}=U^{(x)}_{ai},\qquad
D^{(x)}_{ij}=0,\qquad
D^{(x)}_{ab}=0.
\end{equation}

Похідна фокіана:
\begin{equation}
F^{(y)}=\frac{\partial F}{\partial y}
=
\frac{\partial h}{\partial y}
+
\frac{\partial (J-K)}{\partial y}.
\end{equation}

Неявний внесок:
\begin{equation}\label{eq:HF_Hessian_implicit}
\mathrm{Tr}\!\left[
D^{(x)}\,F^{(y)}
\right]
=
2\sum_{ai}
U^{(x)}_{ai}\,
F^{(y)}_{ai}.
\end{equation}

% ------------------------------------------------------------
\paragraph{Підсумковий вираз ядерного гесіана HF.}
% ------------------------------------------------------------

Збираючи явну і неявну частини, отримуємо:
\begin{equation}\label{eq:HF_Hessian_final}
H_{xy}
=
\mathrm{Tr}\!\left[
\mathbf D\,\frac{\partial^2 \mathbf h}{\partial x\,\partial y}
\right]
+
\frac12
\mathrm{Tr}\!\left[
\mathbf D\,\frac{\partial^2 (\mathbf J-\mathbf K)}{\partial x\,\partial y}
\right]
+
2\sum_{ai}
U^{(x)}_{ai}\,
F^{(y)}_{ai}.
\end{equation}

Перші два доданки — \emph{геометричні}, останній — \emph{орбітальна реакція},
що обчислюється розв’язанням рівнянь \texttt{CPHF}.

%% ------------------------------------------------------------
%\subsection{Структурована процедура обчислення}
%% ------------------------------------------------------------
%
%\begin{enumerate}
%\item Виконати незбурений HF та зберегти $\mathbf C$, $\varepsilon_p$, $(pq|rs)$.
%\item Для кожного збурення $x$ побудувати $g^{(x)}_{ai}$ і розв’язати рівняння \texttt{CPHF}; отримати $U^{(x)}_{ai}$.
%\item Обчислити густину першого порядку $D^{(x)}$.
%\item Обчислити $\partial^2 h/\partial x\,\partial y$ та похідні кулонівської й обмінної частин.
%\item Обчислити похідну фокіана $F^{(y)}$.
%\item Підставити в загальний вираз \eqref{eq:HF_Hessian_final} та отримати елемент гесіана $H_{xy}$.
%\end{enumerate}

